\section{Introduction}

Recent advances in learning-based vulnerability detection have established a
standard methodology. Source code is parsed into a Code Property Graph
(CPG)~\citep{6956589} that integrates syntax, control flow, and data flow. Node
features are extracted using a code-pretrained transformer, and a graph neural
network is subsequently trained to classify functions as vulnerable or not.
Reported performance is consistently strong and has improved steadily across
successive studies. It is also, almost universally, obtained under a random
function-level
split~\citep{chakraborty2020deeplearningbasedvulnerability,ding2024primevul}.

Two questions follow from that, and neither is settled. The first is how much of
the reported performance survives an evaluation that separates what a random
split leaves joined. That evaluation leakage exists is well
established~(Section~\ref{sec:related}); what has not been reported is a
decomposition attributing the loss to each kind of sharing on a single corpus
with the model held fixed. The second is what the code property graph itself
supplies. The evidence offered for the representation is end-to-end comparison
between systems that differ in representation, encoder, training regime and
hyperparameters at once, so a difference between them is not attributable to any
one of those.

We address both by rebuilding the pipeline and varying one factor at a time over
a corpus of 131,903 functions from 949 projects, holding the encoder, depth,
width, pooling, loss, optimizer, splits, folds and metrics constant throughout.
Tightening only what the test partition may share with training lowers ROC-AUC
from $0.796$ to $0.585$. Holding the protocol fixed instead and varying only how
the same graphs are encoded moves them by a comparable amount. Neither the
protocol nor the encoding is consistently reported, and between them they account
for more of the spread in published numbers than the architectural differences
that distinguish published systems.

We organise the study around four questions.

\begin{description}[leftmargin=2.6em,itemsep=2pt]
  \item[\textbf{RQ1}] How much of the performance reported under the standard
  protocol is attributable to what the test partition shares with training,
  and which kind of sharing accounts for how much of it?
  \item[\textbf{RQ2}] Within a fixed pipeline, how much of what a code
  property graph is measured to contribute is set by how it is encoded rather
  than by the representation itself, and which encoding choices carry that
  weight?
  \item[\textbf{RQ3}] What does the representation the network builds from the
  graph capture, and what does it not?
  \item[\textbf{RQ4}] How does the resulting detector compare with a published
  graph-free model~\citep{fu2022linevul} evaluated on identical splits?
\end{description}

Our contributions are as follows:
\begin{enumerate}
    \item \textbf{A decomposition of evaluation leakage.} Across three
    protocols on one corpus with the model held fixed, 47.0\% of the signal
    above chance is attributable to duplicates, patched twins and shared
    commits, 24.4\% to repository overlap, and 28.5\% survives every control.
    The pattern reproduces across four architecture families. Every split we
    release ships with a printed audit of what it separates.

    \item \textbf{The training recipe is a larger factor than any
    representational choice we tested, and is never reported.} Fixing the
    alternative by grid search on the validation partition alone, with the test
    partition withheld until the ranking was settled, the selected recipe is
    worth $+0.032$ ROC-AUC ($p=0.0001$) over the library defaults under
    repository-disjoint evaluation. Whether ensembling helps is itself
    conditional on it: five-fold averaging gains $+0.033$ ($p=0.007$) under the
    default recipe and $-0.002$ ($p=0.59$) under the selected one, so it
    compensates for undertrained members rather than improving the method. The
    extra fifth of training data accounts for $+0.003$ ($p=0.73$). The leakage
    share of Section~\ref{sec:ladder} moves with the same choice, from 71\%
    under the default recipe to 45\% under the selected one: a benchmark does
    not have a leakage figure, a benchmark and a training configuration jointly
    do.

    \item \textbf{Validation is a weak guide to test under a
    repository-disjoint split.} Across the twelve configurations searched,
    validation ROC-AUC spans $0.015$ and test spans $0.029$, with a rank
    correlation of $\rho=0.32$ ($p=0.31$). Selecting honestly on validation
    returns $0.650$ where the grid's best test score is $0.661$, so roughly a
    point of ROC-AUC is unreachable without consulting the test partition. This
    bounds how much of a published figure can be attributed to method rather
    than to search.

    \item \textbf{Component-level conclusions invert with the procedure.}
    Allowing each edge to be traversed in both directions is worth $+0.033$
    ($p=0.012$) under the ensemble and $-0.006$ ($p=0.41$) under the single
    model, on identical data with identical features. Preserving the four
    relation labels is null under both ($+0.002$, 95\% CI $-0.005$ to $+0.010$,
    under the ensemble; $-0.005$, $p=0.22$, under the single model). We therefore report these as
    conditional measurements rather than as recommendations, and argue that
    ablations in this area are uninterpretable unless the training procedure is
    stated with them.

    \item \textbf{A positive control for structural learning.} A null ablation
    is uninformative on its own, since a pipeline that cannot read structure
    and a corpus that contains none produce the same number, and this
    literature reports no control that separates them. Given the same graphs
    with node content replaced by noise, the same models recover labels defined
    by the edge set alone, gaining $+0.104$ and $+0.303$ ROC-AUC from the edges.

    \item \textbf{An evaluation that the aggregate metric hides.} Restricted
    to the $2{,}067$ vulnerable/patched pairs inside the test partition, where
    project identity and style cancel by construction, no configuration tested
    here separates a function from its own fix at better than chance: $0.502$
    on the 165 pairs where only structure differs, and $0.497$ once edges are
    bidirectional.
\end{enumerate}
